In this section, we review the necessary background on Bailey pairs. A Bailey pair relative to~$x$~\cite{Andmultiple} is a pair of sequences $(\alpha_n,\beta_n)$,
$n\in \ZZ_{\geq 0}$,
satisfying
\begin{equation*} %\label{pairdef}
	\beta_n = \sum_{k=0}^n \frac{\alpha_k}{(q)_{n-k}(xq)_{n+k}}.
\end{equation*}
Note that in the limit, this gives (subject to convergence conditions)
\begin{equation} \label{Baileylimit}
	\lim_{n \to \infty} \beta_n = \frac{1}{(q,xq)_{\infty}} \sum_{k=0}^{\infty} \alpha_k.
\end{equation}
In practice, the above sum often becomes an infinite product by an appeal to the Jacobi triple product identity \cite[equation~(2.2.10)]{AndBook},
\begin{align}
	\sum_{n\in\ZZ}(-1)^n z^nq^{n^2}=\bigl(q^2,zq,z^{-1}q;q^2\bigr)_\infty,
	\label{eqn:jtp}
\end{align}
or the quintuple product identity \cite{C-quint},
\begin{align}
	\Biggl(\sum_{\substack{n\in\ZZ\\ n\equiv 0\,\,(\mathrm{mod}\,\,3)}}
	-\sum_{\substack{n\in\ZZ\\ n\equiv 2\,\,(\mathrm{mod}\,\,3)}}
	\Biggr) z^nq^{\frac{1}{3}\binom{n+1}{2}}=\bigl(q,zq,z^{-1}\bigr)_\infty\bigl(z^2q,z^{-2}q;q^2\bigr)_\infty
	\label{eqn:qtp}.
\end{align}

The most important fact about Bailey pairs is the following, which is known as the Bailey lemma. From a given Bailey pair relative to $x$, it produces new Bailey pairs relative to $x$.

\begin{Lemma}[\cite{Andmultiple}] \label{Baileylemma}
	If $(\alpha_n,\beta_n)$ is a Bailey pair relative to $x$, then so is $(\alpha_n',\beta_n')$, where
	\begin{equation*} %\label{Baileylemmaalpha}
		\alpha_n' = \frac{(b,c)_n(xq/bc)^n}{(xq/b,xq/c)_n}\alpha_n
	\qquad
	\text{and}
	\qquad %\label{Baileylemmabeta}
		\beta_n' = \frac{1}{(xq/b,xq/c)_n} \sum_{j=0}^n \frac{(b,c)_j(xq/bc)_{n-j} (xq/bc)^j}{(q)_{n-j}} \beta_j.
	\end{equation*}
\end{Lemma}

A result similar to Lemma \ref{Baileylemma} takes a Bailey pair relative to $x$ and produces Bailey pairs relative to $x/q$.
\begin{Lemma}[\cite{AABLattice}] \label{Baileylatticelemma}
	If $(\alpha_n,\beta_n)$ is a Bailey pair relative to $x$, then $(\alpha_n',\beta_n')$ is a Bailey pair relative to $x/q$, where
	\begin{equation*} %\label{Baileylatticealpha}
		\alpha_n' = (1-x) \Bigl(\frac{x}{bc} \Bigr) \frac{(b,c)_n(x/bc)^n}{(x/b,x/c)_n}\biggl(\frac{\alpha_n}{1-xq^{2n}} - \frac{xq^{2n-2}\alpha_{n-1}}{1-xq^{2n-2}} \biggr)
	\end{equation*}
and
	\begin{equation*} %\label{Baileylatticebeta}
		\beta_n' = \frac{1}{(x/b,x/c)_n} \sum_{j=0}^n \frac{(b,c)_j(x/bc)_{n-j} (x/bc)^j}{(q)_{n-j}} \beta_j.
	\end{equation*}
Here by convention we take
	\begin{equation} \label{eqn:latticeextracond}
		\alpha_{-1} = 0.
	\end{equation}
\end{Lemma}


Finally, we have two lemmas that change a Bailey pair relative to $x$ to one relative to $xq$. The first will be key to introducing $q$-binomial coefficients of the form $\bigl[\begin{smallmatrix} n_{i+1} + 1 \\ n_i \end{smallmatrix}\bigr]$. The second is not needed for the proof of Theorem \ref{CLSXYthm} but will be used to establish a certain Bailey pair later in the paper.
\begin{Lemma}[\cite{Lo1}] \label{littlelemma}
	If $(\alpha_n,\beta_n)$ is a Bailey pair relative to $x$, then $(\alpha_n',\beta_n')$ is a Bailey pair relative to $xq$, where
	\begin{equation*} %\label{littlelemmaalpha}
		\alpha_n' = \frac{1}{1-xq} \biggl(\frac{1-q^{n+1}}{1-xq^{2n+2}} \alpha_{n+1} + \frac{q^n(1-xq^n)}{1-xq^{2n}} \alpha_n \biggr)
	\qquad
	\text{and}
	\qquad %\label{littlelemmabeta}
		\beta_n' = \bigl(1-q^{n+1}\bigr)\beta_{n+1}.
	\end{equation*}
\end{Lemma}
